\clearpage
\section{真人交流参考分与结果复用}
\label{sec:design-favorability}
同一段真人聊天里，甲和乙看到的参考分可以不同。系统按交流方向从聊天历史中统计。例如，甲查看的是乙对自己的回复速度、交流持续性和相对字数。双方都开启后，才能在会话里查看各自方向的结果。

计算时用 $R$ 表示回复速度，$C$ 表示交流持续性，$L$ 表示相对字数，$A$ 表示有效的模型分。具体计算如下。
\begin{equation}
F=\operatorname{clip}_{[0,100]}\left(\begin{cases}
\operatorname{round}[0.25(R+C+L+A)],&\text{有有效模型分},\\
\operatorname{round}[0.35R+0.35C+0.30L],&\text{无有效模型分}.
\end{cases}\right).\label{eq:design-favorability}
\end{equation}
模型分有效时四项等权；没有有效模型分时，按公式中的权重合成三项行为因素。结果取整并限制在0到100之间。例如 $R=90,C=80,L=70,A=84$，综合参考分为81。

\begin{table}[H]
\centering\small
\caption{交流参考的主要使用流程}\label{tab:design-favorability-cache}
\begin{tabularx}{\linewidth}{@{}p{32mm}X@{}}
\toprule 操作阶段 & 处理方式与用户结果 \\ \midrule
开启功能 & 双方开启交流参考后，在会话中查看各自方向的结果 \\
整理交流表现 & 根据聊天历史归纳回应速度、交流延续和表达投入 \\
查看参考结果 & 展示综合分及各项因素，方便了解评分的组成 \\
更新聊天表现 & 随交流推进更新行为因素，保持结果与消息记录衔接 \\
主动分析 & 用户按需补充近期对话的内容分析，取得更完整的交流概括 \\
\bottomrule
\end{tabularx}
\end{table}

查询时以会话编号和分析对象编号定位缓存。缓存可读且末条消息编号一致时，直接返回整份结果。否则根据现有消息重算三项规则指标，旧缓存中可复用的模型分仍可参与计算。查询本身不调用模型。用户主动发起内容分析时才请求模型，模型不可用时以三项规则指标计算参考分。

